\documentclass[11pt]{article}
\title{How the Codynamic Book Machine Works}
\author{Arthur Petron}
\usepackage[margin=1in]{geometry}
\usepackage{graphicx}
\usepackage{amsmath,amssymb,mathtools}
\usepackage{tikz}
\usetikzlibrary{positioning,arrows.meta}
\usepackage{hyperref}
\graphicspath{{media/}{media/diagrams/}{images/}{artwork/}}

\begin{document}

\section*{Abstract}

The Codynamic Book Machine is a schema-backed, message-coordinated, multi-agent authoring system for producing books as living technical objects. Rather than treating a manuscript as a static sequence of files, the system represents authorial intent, outline structure, section drafts, editorial maintenance, references, diagrams, and compiled artifacts as coordinated state surfaces that evolve together through explicit tasks and durable messages. This paper presents the machine as both a writing environment and a working case study in codynamic architecture: a system in which content, code, agents, logs, and outputs remain mutually inspectable and continuously revisable.

The central claim is that long-form technical writing becomes more reliable when the book is modeled as a canonical work object and decomposed into section-level responsibilities handled by specialized agents. In this design, a hypervisor orchestrates work, section agents draft and revise local prose, a gardener agent performs editorial maintenance, a references agent owns citation support, and a diagram agent manages visual assets. These agents do not mutate the manuscript invisibly; they operate through queued tasks, append-only messages, and registered artifacts, so that every substantive change is traceable and reviewable.

The paper focuses on the repository that implements this workflow and uses it as evidence for the architecture it describes. It explains how outline intake becomes normalized structure, how task queues and callbacks coordinate agent work, how citations and diagrams are routed through specialist handoffs, and how generated \LaTeX{} payloads are separated from imported source material. The contribution is therefore twofold: a concrete authoring system that can be inspected and extended, and a technical account of how books can be managed as durable, machine-readable, and collaboratively maintained systems rather than as isolated text files.


\end{document}
